\documentclass[5pt]{article}

\usepackage{sectsty}
\usepackage{graphicx}
\usepackage{lipsum} % for generating dummy text
\usepackage[margin=1in]{geometry}
\usepackage{setspace}
\usepackage{array}
\usepackage{cellspace}
\usepackage{tabularx}
\usepackage{eurosym}
\usepackage[table]{xcolor}
\usepackage{tabularray}
\usepackage{wrapfig}
\usepackage{flafter} 

% package and setup for tables
\usepackage{float}
\usepackage[table]{xcolor}
\renewcommand{\arraystretch}{1.5}
\arrayrulecolor{black}

\usepackage{hyperref}
\usepackage{scrextend}
\graphicspath{ {../assets} }


% Margins
\topmargin=-0.45in
\evensidemargin=0in
\oddsidemargin=0in
\textwidth=6.5in
\textheight=9.0in
\headsep=0.25in

\title{Piano di Progetto}
\author{Jackpot Coding}
\renewcommand*\contentsname{Indice}
\date{\today}

\DefTblrTemplate{contfoot-text}{default}{}
\DefTblrTemplate{conthead-text}{default}{}
\DefTblrTemplate{caption}{default}{}
\DefTblrTemplate{conthead}{default}{}
\DefTblrTemplate{capcont}{default}{}

%STARTOF THE DOCUMENT
\begin{document}

%-------------------------

% Reduce top margin only on the first page
\newgeometry{top=0.5in}

%UNIPD LOGO
  \vspace{8pt}
  \includegraphics[scale=0.2]{UNIPDFull.png}
%END UNIPD LOGO

\vspace{30pt}

%COURSE INFO
\begin{minipage}[t]{0.48\textwidth}
    %COURSE TITLE
        \begin{flushleft}
            Informatica\\
            \vspace{5pt}
            \textbf{\LARGE Ingegneria del Software}\\
            Anno Accademico: 2023/2024
        \end{flushleft}
    %END COURSE TITLE
\end{minipage}
%END COURSE INFO


\vspace{5px}


%BLACK LINE
\rule{\textwidth}{5pt}

%JACKPOT CODING INFO
\begin{minipage}[t]{0.50\textwidth}
    %LOGO JACKPOT CODING
    \begin{flushleft}
        \hspace{10pt}
        \includegraphics[scale=0.65]{jackpot-logo.png} 
    \end{flushleft}
\end{minipage}
\hspace{-60pt} % This adds horizontal space between the minipages
\begin{flushright}
    \begin{minipage}[t]{0.50\textwidth}
        %INFO JACKPOT CODING
        \begin{flushright}
            Gruppo: {\Large Jackpot Coding}\\
            Email: \href{mailto:jackpotcoding@gmail.com}{jackpotcoding@gmail.com}
        \end{flushright}
    \end{minipage}
\end{flushright}
%END JACKPOT CODING INFO

\vspace{24pt}

%TITLE
\begin{center}
    \textbf{\LARGE PIANO DI PROGETTO}
\end{center}
%END TITLE

\vspace{13pt}

\begin{flushleft}
    \begin{spacing}{1.5}
        REDATTORE: M. Gobbo, M. Favaretto\\%INSERT HERE THE NAMES
        VERIFICATORI: G. Moretto, E. Gallo, M.Gobbo\\
        \vspace{7pt}
        DESTINATARI: Prof. T. Vardanega, Prof. R. Cardin\\%INSERT HERE THE NAMES
    \end{spacing}
\end{flushleft}

\begin{flushright}
    \begin{spacing}{1}
        USO: ESTERNO\\
        VERSIONE: 1.0.2\\
    \end{spacing}
\end{flushright}


% Restore original margins from the second page onwards
\restoregeometry

\pagebreak

\textbf{\Large Registro delle modifiche}
\begin{longtblr}
	{
		colspec={|Q[0.10\linewidth]|Q[0.10\linewidth]|Q[0.15\linewidth]|Q[0.15\linewidth]|Q[0.45\linewidth]|},
		rows={halign=l},
		column{1}={halign=c},
		column{3}={halign=c},
		column{4}={halign=c},
		row{1}={halign=c},
		row{odd} = {gray!20},
		row{1}={teal!50},
	}
	\hline
    \textbf{Versione} & \textbf{Data} & \textbf{Autore} & \textbf{Verificatore} & \textbf{Modifica} \\
    \hline
    v1.0.2 & 04/04/2024 & G. Moretto & - & Rimozione Redattori e Verificatori come da feedback RTB\\
    \hline
    v1.0.1 & 22/03/2024 & G. Moretto & M.Gobbo & Aggiornamento tabella consuntivo RTB \\
    \hline
    v1.0.0 & 18/03/2024 & - & E. Gallo & Approvazione documento \\
    \hline
    v0.0.17 & 15/03/2024 & M. Gobbo - M. Favaretto & G. Moretto & Scrittura sezione 6 \\
    \hline
    v0.0.16 & 14/03/2024 & M. Gobbo - M. Favaretto & G. Moretto & Fine sezione 5 \\
    \hline
    v0.0.15 & 13/03/2024 & M. Gobbo - M. Favaretto & G. Moretto & Modifiche sezione 5 \\
    \hline
    v0.0.14 & 11/03/2024 & M. Gobbo - M. Favaretto & G. Moretto & Modifiche sezione 5 \\
    \hline
    v0.0.13 & 05/03/2024 & M. Gobbo - M. Favaretto & G. Moretto & Modifiche sezione 4 \\
    \hline
    v0.0.12 & 16/02/2024 & M. Gobbo - M. Favaretto & G. Moretto & Modifiche sezione 4 \\
    \hline
    v0.0.11 & 07/02/2024 & M. Gobbo - M. Favaretto & G. Moretto & Scrittura sezione 5 \\
    \hline
    v0.0.10 & 06/02/2024 & M. Gobbo - M. Favaretto & G. Moretto & Scrittura sezione 5 \\
    \hline
    v0.0.9 & 05/02/2024 & M. Gobbo - M. Favaretto & G. Moretto & Scrittura sezione 5 \\
    \hline
    v0.0.8 & 12/01/2024 & M. Gobbo - M. Favaretto & G. Moretto & Migliore indentazione del documento \\
    \hline
    v0.0.7 & 12/12/2023 & M. Gobbo - M. Favaretto & G. Moretto & Completamento sezione 4 \\
    \hline
    v0.0.6 & 23/12/2023 & M. Gobbo - M. Favaretto & G. Moretto & Scrittura sezione 4 \\
    \hline
    v0.0.5 & 20/11/2023 & M. Gobbo - M. Favaretto &G. Moretto- & Scrittura sezione 4 \\
    \hline
    v0.0.4 & 17/11/2023 & M. Gobbo - M. Favaretto & G. Moretto & Scrittura sezione 3 \\
    \hline
    v0.0.3 & 16/11/2023 & M. Gobbo - M. Favaretto & G. Moretto & Completamento sezione 2 \\
    \hline
    v0.0.3 & 15/11/2023 & M. Gobbo - M. Favaretto & G. Moretto & Scrittura sezione 2 \\
    \hline
    v0.0.2 & 10/11/2023 & M. Gobbo - M. Favaretto & G. Moretto & Scrittura sezione 1 \\
    \hline
    v0.0.1 & 03/11/2023 & M. Gobbo - M. Favaretto & G. Moretto & Creata struttura del documento \\
    \hline
\end{longtblr}

\pagebreak
\tableofcontents
\pagebreak

\begin{flushleft}
	\textbf{\Large Elenco delle tabelle}
	\begin{spacing}{1.5}
		Tabella 1: Preventivo prima parte RTB \\
		Tabella 2: Preventivo seconda parte RTB \\
		Tabella 3: Preventivo rifiniture RTB \\
		Tabella 4: Piano ore preventivate per il primo periodo \\
		Tabella 5: Costi preventivati primo periodo \\
		Tabella 6: Preventivo fase preliminare PB \\
		Tabella 7: Preventivo progettazione e programmazione PB \\
		Tabella 8: Preventivo progettazione secondaria PB \\
		Tabella 9: Preventivo fase finale PB \\
		Tabella 10: Piano ore preventivate per il secondo periodo \\
		Tabella 11: Costi preventivati primo periodo \\
		Tabella 12: Resoconto ore e costi RTB \\
	\end{spacing}
\end{flushleft}

\pagebreak

\section{Introduzione}
  \subsection{Scopo del documento}
    Il presente documento ha come fine la presentazione della pianificazione e delle modalità di sviluppo di questo progetto. 
    Nella fattispecie verranno esposte le analisi dei possibili rischi con alcune proposte di mitigazione, il modello adottato dal \textit{team}, 
    il preventivo e il consuntivo di periodo.

  \subsection{Scopo del capitolato}
    Le \textit{I.A.\textsuperscript{G}} stanno vivendo un momento di grande innovazione ed entusiasmo, 
    riuscendo a cogliere l'attenzione anche di utenti al di fuori dell'ambito informatico lavorativo e accademico, 
    si veda come \textit{ChatGPT\textsuperscript{G}} sia diventato un fenomeno culturale. \\
    La versatilità dell'\textit{I.A.\textsuperscript{G}} è oggi al centro dell'attenzione di molte \textit{software house\textsuperscript{G}}, 
    poiché posso essere usate anche per migliorare e velocizzare la produzione. 
    Per farne buon uso, è però necessario essere in grado di fornire i corretti \textit{prompt\textsuperscript{G}} al 
    modello di \textit{I.A.\textsuperscript{G}} in uso. \\


    L'obiettivo di questo progetto è realizzare un \textit{software\textsuperscript{G}} in grado di generare 
    un \textit{prompt\textsuperscript{G}} a partire da una richiesta in linguaggio naturale\textsuperscript{G}. La richiesta dovrà riguardare un'interrogazione 
    di un \textit{database\textsuperscript{G}} caricato sul sistema dall'utente. Tale \textit{prompt\textsuperscript{G}} sarà successivamente da fornire 
    ad un \textit{LLM\textsuperscript{G}}, il quale restituirà all'utente una \textit{query\textsuperscript{G}} nel linguaggio 
    \textit{SQL\textsuperscript{G}}.

  \subsection{Glossario}
    Alcuni termini presenti in questo documento potrebbero generare incomprensioni o necessitare di chiarimenti. 
    Al fine di evitare queste eventualità, tali termini sono contrassegnati dalla lettera \textit{G} maiuscola posta ad apice della parola, 
    per indicare che la loro spiegazione è presente all'interno del documento \textit{Glossario}.

  \subsection{Riferimenti}
    \subsubsection{Riferimenti Informativi}
      Slide del corso "Inegegneria Del Software"
      \begin{itemize}
            \item \href{https://www.math.unipd.it/~tullio/IS-1/2023/Dispense/T3.pdf}{Modelli di sviluppo software}
            \item \href{https://www.math.unipd.it/~tullio/IS-1/2023/Dispense/T4.pdf}{Gestione del progetto}
            \item \href{https://www.math.unipd.it/~tullio/IS-1/2023/Dispense/T5.pdf}{Analisi dei requisiti} 
      \end{itemize}
    \subsubsection{Riferimenti Normativi}
      \begin{itemize}
            \item Documento "Norme di progetto"
            \item \href{https://www.math.unipd.it/~tullio/IS-1/2023/Dispense/PD2.pdf}{Regolamento del progetto}
            \item \href{https://www.math.unipd.it/~tullio/IS-1/2023/Progetto/C9.pdf}{Capitolato C9 "ChatSQL"} 
      \end{itemize}

\section{Analisi dei rischi}
  La realizzazione di un progetto \textit{software\textsuperscript{G}} comporta una serie di rischi che devono essere valutati e gestiti attentamente al fine di garantire il successo del prodotto finale. Nel corso di questo paragrafo, esamineremo i rischi che siamo riusciti ad individuare e le strategie per mitigarli efficacemente. \\[1\baselineskip]
  Per facilitare l'analisi, ogni rischio avrà la seguente struttura:
  \begin{itemize}
    \item \textbf{Codice}: Identifica univocamente il rischio, che può essere di 3 tipi:
        \begin{itemize}
        \item \textbf{RT}: Rischi tecnologici
        \item \textbf{RO}: Rischi organizzativi
        \item \textbf{RC}: Rischi comunicativi
        \end{itemize}
    \item \textbf{Titolo}: Breve titolo descrittivo del rischio.
    \item \textbf{Descrizione}: Spiegazione dettagliata del rischio.
    \item \textbf{Occorrenza}: Frequenza con cui il rischio si verifica.
    \item \textbf{Pericolosità}: Livello di pericolo associato al rischio.
    \item \textbf{Mitigazione}: Azioni per ridurre l'impatto o la probabilità del rischio.
  \end{itemize}

  \subsection{Rischi tecnologici}
    \textbf{RT-1} \\
    \textbf{Inesperienza nello sviluppo tecnologico}
    \begin{itemize}
      \item \textbf{Descrizione}: La mancanza di esperienza nel gruppo su una determinata tecnologia o piattaforma potrebbe rallentare lo sviluppo del \textit{software\textsuperscript{G}} e portare a errori nella progettazione e nell'implementazione\textsuperscript{G}.
      \item \textbf{Occorrenza}: Media
      \item \textbf{Pericolosità}: Media
      \item \textbf{Mitigazione}: 
        \begin{itemize}
          \item Assicurarsi che i membri del gruppo ricevano formazione e supporto adeguati sulla tecnologia o piattaforma specifica.
          \item Affidare compiti critici a membri del gruppo con esperienza nel settore.
          \item Effettuare una pianificazione accurata per consentire più tempo per l'apprendimento e l'adattamento.
          \end{itemize}
    \end{itemize}
    \textbf{RT-2} \\
    \textbf{Inesperienza strumenti \textit{software}}
    \begin{itemize}
      \item \textbf{Descrizione}: Il gruppo non ha familiarità con l'utilizzo di \textit{software\textsuperscript{G}} specifici per la gestione di un progetto. Questa mancanza di esperienza potrebbe rallentare la produttività del gruppo e compromettere la qualità del progetto.
      \item \textbf{Occorrenza}: Bassa
      \item \textbf{Pericolosità}: Media
      \item \textbf{Mitigazione}: 
        \begin{itemize}
          \item  Attivarsi per avere una formazione adeguata sull'uso del \textit{software\textsuperscript{G}} gestionali attraverso \textit{tutorial\textsuperscript{G}} e documentazione.
          \item Controllare se gli strumenti scelti dispongono di una solida guida utente e documentazione varia.
        \end{itemize}
    \end{itemize}

  \subsection{Rischi interni}
    \subsubsection{Rischi organizzativi}
      \textbf{RO-1}\\
      \textbf{Organizzazione carente}
      \begin{itemize}
        \item \textbf{Descrizione}: Una struttura organizzativa inefficiente o una mancanza di chiarezza nei ruoli e nelle responsabilità possono portare a ritardi nelle decisioni, duplicazioni di compiti e confusione all'interno del gruppo.
        \item \textbf{Occorrenza}: Alta
        \item \textbf{Pericolosità}: Media
        \item \textbf{Mitigazione}: 
          \begin{itemize}
            \item Definire chiaramente le responsabilità di ciascun membro del gruppo e comunicarle in modo trasparente a tutti.
            \item Condurre regolarmente riunioni di coordinamento per monitorare lo stato del progetto e risolvere eventuali problemi organizzativi in modo tempestivo.
          \end{itemize}
      \end{itemize}
      \textbf{RO-2}\\
      \textbf{Distribuzione disomogenea dei compiti}
      \begin{itemize}
        \item \textbf{Descrizione}: Una distribuzione disomogenea dei compiti all'interno del gruppo potrebbe portare a una mancanza di lavoro per alcuni componenti e sovraccarico di lavoro per altri, compromettendo l'efficienza\textsuperscript{G} complessiva del progetto.
        \item \textbf{Occorrenza}: Media
        \item \textbf{Pericolosità}: Alta
        \item \textbf{Mitigazione}: 
          \begin{itemize}
            \item Pianificare attentamente l'allocazione dei compiti, tenendo conto delle competenze, delle disponibilità di ciascun membro del gruppo.
            \item Promuovere una cultura di collaborazione e flessibilità, incoraggiando il supporto reciproco tra i membri del gruppo e la condivisione dei carichi di lavoro.
            \item Monitorare regolarmente lo stato dei compiti assegnati e intervenire prontamente per ridistribuire le risorse in caso di necessità.
          \end{itemize}
      \end{itemize}
    \subsubsection{Rischi comunicativi}
      \textbf{RC-1}\\
      \textbf{Mancanza di comunicazione interna efficace}
      \begin{itemize}
        \item \textbf{Descrizione}: La mancanza di una comunicazione interna efficace tra i membri del gruppo può portare a fraintendimenti, duplicazione del lavoro e rallentamenti nel processo decisionale.
        \item \textbf{Occorrenza}: Bassa
        \item \textbf{Pericolosità}: Alta
        \item \textbf{Mitigazione}: 
          \begin{itemize}
            \item Stabilire procedure chiare e canali di comunicazione ben definiti all'interno del gruppo.
            \item Organizzare riunioni regolari per discutere lo stato del progetto, assegnare compiti e condividere aggiornamenti.
            \item Utilizzare strumenti di gestione del progetto e di collaborazione per facilitare la comunicazione e 
            la condivisione delle informazioni come \textit{Trello} (utilizzato come \textit{dashboard\textsuperscript{G}}).
          \end{itemize}
      \end{itemize}
      \textbf{RC-2}\\
      \textbf{Problemi di comunicazione con l'azienda}
      \begin{itemize}
        \item \textbf{Descrizione}: La comunicazione inefficace con l'azienda può portare a malintesi, ritardi nelle consegne e insoddisfazione del cliente, rallentando inoltre l'attività del gruppo.
        \item \textbf{Occorrenza}: Bassa
        \item \textbf{Pericolosità}: Alta
        \item \textbf{Mitigazione}: 
          \begin{itemize}
            \item Stabilire canali di comunicazione chiari e definire ruoli e responsabilità per l'interazione con l'azienda.
            \item Fornire aggiornamenti regolari e trasparenti sullo stato del progetto.
            \item Ascoltare attivamente i \textit{feedback\textsuperscript{G}} dell'azienda e risolvere eventuali problemi di comunicazione tempestivamente.
          \end{itemize}
      \end{itemize}
      \textbf{RC-3}\\
      \textbf{Tensioni e conflitti interni}
      \begin{itemize}
        \item \textbf{Descrizione}: La presenza di tensioni, conflitti o mancanza di collaborazione tra i membri del gruppo può compromettere l'efficienza e la coesione del gruppo, influenzando negativamente la qualità del lavoro e il raggiungimento degli obiettivi del progetto.
        \item \textbf{Occorrenza}: Media
        \item \textbf{Pericolosità}: Media
        \item \textbf{Mitigazione}: 
          \begin{itemize}
            \item Utilizzare la votazione per prendere decisioni oggettive basate sulla maggioranza delle preferenze. Assicurarsi che tutte le parti coinvolte comprendano e accettino il risultato della votazione per favorire la collaborazione e il senso di appartenenza al gruppo.
            \item Stabilire procedure per la gestione dei conflitti e incoraggiare il coinvolgimento di un mediatore neutrale (arbitro) se necessario.
            \item Utilizzare il metodo del dibattito per esaminare i punti di vista contrastanti in modo strutturato e razionale
          \end{itemize}
      \end{itemize}

\section{Modello di Sviluppo}
  \subsection{Modello \textit{Agile}}
    Il gruppo \textit{"Jackpot Coding"} ha deciso di ispirarsi ai Modelli Agili per lo sviluppo del progetto.\\
    Il modello\textsuperscript{G} è caratterizzato da rilasci continui e un progressivo arricchimento delle funzionalità del prodotto, questo approccio consente di identificare agevolmente ciascun requisito\textsuperscript{G}, assegnandogli una priorità e organizzando lo sviluppo in modo da mantenere il prodotto funzionante in ogni fase, attivamente perseguendo i requisiti e adattandosi lungo il percorso verso il completamento del prodotto finale.\\
    Anche la documentazione creata durante le varie fasi di sviluppo sarà costantemente aggiornata per integrare e/o modificare le informazioni rilevanti.

  \subsection{Motivazione}
    I motivi che hanno spinto il gruppo ad adottare un modello\textsuperscript{G} di questo tipo sono molteplici e sono i seguenti:
    \begin{itemize}
        \item Attraverso la suddivisione temporale delle attività, i \textit{test} diventano più agevoli, semplificando la fase di modifica e \textit{test}.
        \item Gli errori nel prodotto sono più facilmente rilevabili grazie ai continui incrementi\textsuperscript{G}.
        \item La capacità di adattamento del \textit{team} alle variazioni è migliorata, con una documentazione che si evolve in modo collaborativo e flessibile insieme al prodotto, adattandosi alle esigenze del progetto.
        \item Le sfide nelle varie fasi vengono individuate e possono essere affrontate tramite la modifica degli obiettivi, se necessario, per risolvere eventuali problemi di natura temporale o organizzativa.
    \end{itemize}

%-----------------------------

\section{Pianificazione}
  \textbf{Scadenze}:
  \begin{itemize}
      \item \textit{Requirments and Technology Baseline\textsuperscript{G}}: 01/03/2024
      \item \textit{Product Baseline\textsuperscript{G}}: 17/05/2024
  \end{itemize}
  % --- Aggiunta possibile linea del tempo

  \subsection{Verso la \textit{RTB}}
    \textbf{Periodo:} dal 24/10/2023 al 24/03/2024
    \subsubsection{Preparazione preliminare}
      \textbf{Periodo:} dal 24/10/2023 al 15/11/2023 
      \vspace{0.3cm} \\
      In questa fase il \textit{team} appena formato deve concentrarsi nella scelta del progetto dopo le necessarie consulenze con i proponenti presi in considerazione. \\
      Contemporaneamente verranno scelti alcuni dei principali sistemi di gestione del progetto e della comunicazione. \\
      Nel nostro caso, tali scelte sono ricadute nelle seguenti tecnologie: 
      \begin{itemize}
        \item \textit{Trello} per la gestione di \textit{issue\textsuperscript{G}} e \textit{ticketing\textsuperscript{G}} delle attività
        \item \textit{Telegram\textsuperscript{G}} e \textit{Discord\textsuperscript{G}} per la comunicazione asincrona\textsuperscript{G} e sincrona\textsuperscript{G}
        \item \textit{\LaTeX\textsuperscript{G}} come linguaggio per la stesura dei documenti
      \end{itemize}
      È necessario iniziare una parziale identificazione di alcuni rischi, riportati nell'apposita sezione, che potrebbero compromettere il corretto svolgimento del progetto, 
      ma utili anche a presentare un preventivo maggiormente accurato. 
      In ultimo, è importante fornire almeno un'iniziale stesura del documento di \textit{Way of Working\textsuperscript{G}}, necessario per definire i principali metodi di lavoro dei membri del \textit{team}.
      % 31/10 -> scelta proposta 
      % 06/11 -> aggiudicazione appalti
      % 15/11 -> verbale ultimato wow
    
    \subsubsection{Analisi Requisiti}
      \textbf{Periodo:} dal 15/11/2023 al 25/01/2024 
      \vspace{0.3cm} \\
      Una volta che l'appalto è stato aggiudicato al \textit{team}, è di principale importanza l'analisi di ogni possibile caso d'uso\textsuperscript{G} e requisito\textsuperscript{G} del capitolato\textsuperscript{G} preso a carico.
      Per evitare dubbi ed avere un \textit{feedback\textsuperscript{G}} sul lavoro svolto, è utile organizzare un colloquio con il proponente relativamente all'analisi svoltasi fino a quel momento. \\
      Dopo una fase di \textit{brainstorming\textsuperscript{G}} del gruppo e \textit{feedback\textsuperscript{G}} dal proponente riguardo i casi d'uso identificati, si procederà alla stesura del documento di Analisi dei Requisiti\textsuperscript{G},
      dove verranno riportati tutti i casi d'uso identificati, i requisiti obbligatori, desiderabili e opzionali. \\
      Questa fase è dunque di fondamentale importanza per la successiva.
      % 15/11 -> inizio discussione adr, contattato proponente
      % ??/01 -> fine adr?
    
    \subsubsection{Progettazione per RTB}
      \textbf{Periodo} dal 18/01/2024 al 12/02/2024 
      \vspace{0.3cm} \\
      % periodo 18-25 ultimazione adr e relativa verifica -> quindi inizio poc
      In questo periodo, basandosi sulla precedente analisi dei requisiti\textsuperscript{G}, inizierà la realizzazione del \textit{PoC\textsuperscript{G}}. \\
      Questa fase si concluderà con un colloquio con il proponente, per un \textit{feedback\textsuperscript{G}} circa il \textit{PoC\textsuperscript{G}} prodotto.
      % 18/01 -> inizio POC
      % 11/02 -> fine POC
    
    \subsubsection{Rifiniture}
      \textbf{Periodo:} dal 12/02/2024 al 24/03/2024 
      \vspace{0.3cm} \\
      Nell'ultima fase della prima parte del progetto, verranno ultimate le modifiche e relative verifiche ai vari documenti 
      prodotti fino a questo momento. \\
      Questa parte si concluderà con un colloquio e una presentazione del lavoro fino ad ora svolto 
      ai professori del corso Ingegneria del \textit{Software}.
      % aggiunta immagini / schemi esempio: linee del tempo
      % 25/02 -> fine rtb (non definitiva)
      % rifinitura ultimi documenti: pdp, pdq, ndp

  %-------------
  
  \subsection{Verso la \textit{PB}}
    \textbf{Periodo Pianificato:} dal 17/03/2024 al 17/05/2024
    \subsubsection{Preparazione preliminare}
      \textbf{Periodo:} dal 17/03/2024 al 24/03/2024 
      \vspace{0.3cm} \\
      In questo periodo l'obiettivo sarà l'integrazione dei \textit{feedback\textsuperscript{G}} e la sistemazione degli errori emersi durante la \textit{RTB\textsuperscript{G}}; 
      una volta passata questa fase preliminare è il momento di procedere verso la \textit{PB\textsuperscript{G}}.
      Verranno dunque pianificate le attività future e la relativa suddivisione dei compiti.
    
    \subsubsection{Progettazione primaria e codifica dei requisiti obbligatori}
      \textbf{Periodo:} dal 24/03/2024 al 15/04/2024 
      \vspace{0.3cm} \\
      Terminata la retrospettiva dell'\textit{RTB\textsuperscript{G}}, Verrà utilizzato il \textit{PoC\textsuperscript{G}}, realizzato nelle fasi precedenti, 
      come riferimento per la costruzione del progetto effettivo.
      Lo scopo di questa fase è la realizzazione di una prima versione del prodotto finale, che dovrà soddisfare tutti i requisiti\textsuperscript{G} obbligatori. \\
      Per raggiungere questo obiettivo, andranno svolte le seguenti attività:
      \begin{itemize}
        \item \textbf{Studio tecnologie}: al fine di produrre un risultato migliore, andranno approfondite maggiormente le tecnologie approcciate 
        durante la codifica del \textit{PoC\textsuperscript{G}}
        \item \textbf{Progettazione}: scelta dei \textit{design pattern\textsuperscript{G}} e delle principali unità architetturali\textsuperscript{G}
        \item \textbf{Realizzazione del prodotto}: fase che si dividerà nelle seguente sotto-attività che si svolgeranno parallelamente
        \begin{itemize}
          \item \textbf{Codifica}: scrittura del codice che soddisferà i requisiti\textsuperscript{G} obbligatori, rispettando le scelte architetturali
          \item \textbf{Documentazione}: ampliamento della documentazione già prodotta e scrittura della documentazione relativa al programma
          \item \textbf{Verifica e validazione}: fase di svolgimento di test automatici e non al fine di verificare che il codice non produca errori e soddisfi i requisiti\textsuperscript{G} obbligatori.
        \end{itemize}  
      \end{itemize}
    
    \subsubsection{Progettazione secondaria e codifica dei requisiti opzionali}
      \textbf{Periodo:} dal 15/04/2024 al 10/05/2024 
      \vspace{0.3cm} \\
      Questo periodo sarà simile al precedente, con la differenza che sarà dedicato alla realizzazione e codifica dei requisiti desiderabili e opzionali 
      che si è deciso di implementare. \\
      Anche la suddivisione di attività risulterà molto simile:
      \begin{itemize}
        \item \textbf{Scelta dei requisiti}: scelta dei requisiti non obbligatori da implementare. Da concordare con i membri del gruppo in base alle risorse attualmente disponibili.
        \item \textbf{Realizzazione del prodotto}: fase che verrà divisa nelle seguente sotto-attività e che sarà incrementale al periodo precedente
        \begin{itemize}
          \item \textbf{Codifica}: scrittura del codice che soddisferà i requisiti non obbligatori, rispettando le scelte architetturali
          \item \textbf{Documentazione}: ampliamento della documentazione già prodotta
          \item \textbf{Verifica e validazione}: fase di svolgimento di test automatici e non al fine di verificare che il codice non produca errori e soddisfi i requisiti non obbligatori integrati.
        \end{itemize}  
      \end{itemize}
    
    \subsubsection{Validazione finale, collaudo e \textit{PB}}
      \textbf{Periodo:} dal 10/05/2024 al 17/05/2024 
      \vspace{0.3cm} \\
      In quest'ultima fase, il \textit{team} ultimerà le attività di verifica, validazione e collaudo del codice prodotto e della documentazione scritta fino a questo momento.
      In caso di necessità, verranno aggiunte le necessarie correzioni risultate dai \textit{feedback\textsuperscript{G}} di queste attività. \\
      Il fine di questo periodo è la seconda revisione del progetto, ossia la \textit{Product Baseline\textsuperscript{G}}.

\section{Preventivo}
  Di seguito sono presentate le tabelle con le ore e i costi preventivati per i periodi \textit{RTB} e \textit{PB} e i loro relativi sottoperiodi.
  Questa pianificazione rispetta quanto è stato dichiarto nel \href{https://jackpot-coding.github.io/chatSQL/docs/esterni/documento-impegni-v.1.1.pdf}{documento di impegni} 
  presentato ad inizio progetto.

  \subsection{Periodo \textit{RTB}}
    \subsubsection{Prima Parte}
      \textbf{Periodo:} dal 24/10/2023 al 25/01/2024
      \begin{table}[H]
     	  \centering
     	  \rowcolors{2}{black!15}{}
     	  \resizebox{\linewidth}{!}{
     		  \begin{tabular}{|c|c|c|c|c|c|c|c|}
     			  \rowcolor{teal!50}
     			  \hline
     			  \textbf{MEMBRI} & \textbf{Responsabile} & \textbf{Amministratore} & \textbf{Analista} & \textbf{Progettista} &	\textbf{Programmatore} & \textbf{Verificatore} & \textbf{Totale} \\
     			  \hline
            Giulio Moretto & 1 & 1 & 9 & - & - & - & 11 \\ \hline
            Riccardo Simionato & 1 & - & 9 & - & - & 8 & 18 \\ \hline
            Marco Favaretto & 3 & 1 & 9 & - & - & 7 & 20 \\ \hline
            Marco Gobbo & 3 & 2 & 9 & - & - & 6 & 20 \\ \hline
            Matteo Camillo & - & 2 & 9 & - & - & 2 & 13 \\ \hline
            Edoardo	Gallo & - & 2 & 9 & - & - & 2 & 13 \\ \hline
          \end{tabular}
        }
        \caption{Tabella mostrante le ore preventivate per la preparazione preliminare e l'analisi dei requisiti}
      \end{table}

    \subsubsection{Seconda Parte}
      \textbf{Periodo} dal 18/01/2024 al 12/02/2024 
      \begin{table}[H]
     	  \centering
     	  \rowcolors{2}{black!15}{}
     	  \resizebox{\linewidth}{!}{
     		  \begin{tabular}{|c|c|c|c|c|c|c|c|}
     		  	\rowcolor{teal!50}
     		  	\hline
     		  	\textbf{MEMBRI} & \textbf{Responsabile} & \textbf{Amministratore} & \textbf{Analista} & \textbf{Progettista} &	\textbf{Programmatore} & \textbf{Verificatore} & \textbf{Totale} \\
     		  	\hline
            Giulio Moretto & 1 & 1 & 2 & 15 & 7 & 1 & 27 \\ \hline 
            Riccardo Simionato & 1 & - & 2 & 6 & - & 4 & 13 \\ \hline
            Marco Favaretto & 3 & 1 & 2 & - & - & 3 & 9 \\ \hline
            Marco Gobbo & 2 & 2 & 2 & - & - & 3 & 9 \\ \hline
            Matteo Camillo & - & 3 & 2 & 9 & 7 & 2 & 23 \\ \hline
            Edoardo	Gallo & - & 2 & 2 & 10 & 7 & 2 & 23 \\ \hline
          \end{tabular}
        }
        \caption{Tabella mostrante le ore preventivate per il periodo precedente al primo colloquio per l'\textit{RTB}}
      \end{table}
    
    \subsubsection{Rifiniture}
      \textbf{Periodo:} dal 12/02/2024 al 24/03/2024 
      \begin{table}[H]
    	  \centering
    	  \rowcolors{2}{black!15}{}
    	  \resizebox{\linewidth}{!}{
    		  \begin{tabular}{|c|c|c|c|c|c|c|c|}
    		  	\rowcolor{teal!50}
    		  	\hline
    		  	\textbf{MEMBRI} & \textbf{Responsabile} & \textbf{Amministratore} & \textbf{Analista} & \textbf{Progettista} &	\textbf{Programmatore} & \textbf{Verificatore} & \textbf{Totale} \\
    		  	\hline
            Giulio Moretto & 1 & 1 & 1 & - & - & - & 3 \\ \hline
            Riccardo Simionato & 1 & - & 1 & - & - & 8 & 10 \\ \hline
            Marco Favaretto & 4 & - & 1 & - & - & 7 & 12 \\ \hline
            Marco Gobbo & 3 & 2 & 1 & - & - & 6 & 12 \\ \hline
            Matteo Camillo & - & 2 & 1 & - & - & 2 & 5 \\ \hline
            Edoardo	Gallo & - & 2 & 1 & - & - & 2 & 5 \\ \hline
          \end{tabular}
        }
        \caption{Tabella mostrante le ore preventivate per le correzioni e le rifiniture postume al primo colloquio per l'\textit{RTB}}
      \end{table}
    %--------------

    \subsubsection{\textit{RTB}}
      Qui sono presentate le ore e i costi preventivati per il periodo \textit{RTB} complessivo.
      \begin{table}[H]
     	  \centering
     	  \rowcolors{2}{black!15}{}
     	  \resizebox{\linewidth}{!}{
     	    \begin{tabular}{|c|c|c|c|c|c|c|c|}
     	      \rowcolor{teal!50}
     	      \hline
     	      \textbf{MEMBRI} & \textbf{Responsabile} & \textbf{Amministratore} & \textbf{Analista} & \textbf{Progettista} &	\textbf{Programmatore} & \textbf{Verificatore} & \textbf{Totale} \\
     	      \hline
            Giulio Moretto & 3 & 3 & 12	& 15 & 7 & 1 & 41 \\ \hline
            Riccardo Simionato & 3 & - & 12	& 5	& - & 21 & 41 \\ \hline
            Marco Favaretto & 10 & 2 & 12 & - & - & 17 & 41 \\ \hline
            Marco Gobbo & 8 & 6 & 12 & - & - & 15 & 41 \\ \hline
            Matteo Camillo & - & 7 & 12 & 10 & 7 & 5 & 41 \\ \hline
            Edoardo	Gallo & - & 6 & 12 & 10 & 7 & 6 & 41 \\ \hline
          \end{tabular}
        }
        \caption{Piano ore preventivate per il primo periodo}
      \end{table}

      \begin{table}[H]
      	\centering
      	\rowcolors{2}{black!15}{}
      	\resizebox{\linewidth}{!}{
      		\begin{tabular}{|c|c|c|c|c|c|c|c|}
      			\rowcolor{teal!50}
      			\hline
      			\textbf{MEMBRI} & \textbf{Responsabile} & \textbf{Amministratore} & \textbf{Analista} & \textbf{Progettista} &	\textbf{Programmatore} & \textbf{Verificatore} & \textbf{Totale} \\
      			\hline
            Totale per ruolo & 24 & 24 & 72 & 40 & 21 & 65 & 246 \\ \hline
            Costo per ora & € 30,00 & € 20,00 & € 25,00 & € 25,00 & € 15,00 & € 15,00 & - \\ \hline 
            Costo totale per ruolo & € 720,00 & € 480,00 & € 1.800,00 & € 1.000,00 & € 315,00 & € 975,00 & € 5.290,00 \\ \hline
          \end{tabular}
        }
        \caption{Tabella costi preventivati per il primo periodo}
      \end{table}

  \subsection{Periodo \textit{PB}}
    \subsubsection{Fase preliminare}
      \textbf{Periodo:} dal 24/03/2024 al 15/04/2024 
      \begin{table}[H]
        \centering
        \rowcolors{2}{black!15}{}
        \resizebox{\linewidth}{!}{
          \begin{tabular}{|c|c|c|c|c|c|c|c|}
            \hline
            \rowcolor{teal!50}
            \textbf{MEMBRI} & \textbf{Responsabile} & \textbf{Amministratore} & \textbf{Analista} & \textbf{Progettista} & \textbf{Programmatore} & \textbf{Verificatore} & \textbf{Totale} \\ \hline
            Giulio Moretto & 1 & 2 & - & - & - & 3 & 6 \\ \hline
            Riccardo Simionato & 1 & 2 & - & - & - & 3 & 6 \\ \hline
            Marco Favaretto & - & - & - & - & - & 3 & 3 \\ \hline
            Marco Gobbo & - & - & - & - & - & 4 & 4 \\ \hline
            Matteo Camillo & 2 & 1 & - & - & - & 2 & 5 \\ \hline
            Edoardo Gallo & 2 & 1 & - & - & - & 3 & 6 \\ \hline
          \end{tabular}
        }
        \caption{Ore di lavoro preventivate per la parte preliminare della \textit{PB}}
      \end{table}
    
    \subsubsection{Progettazione primaria e codifica dei requisiti obbligatori}
      \textbf{Periodo:} dal 24/03/2024 al 15/04/2024 
      \begin{table}[H]
        \centering
        \rowcolors{2}{black!15}{}
        \resizebox{\linewidth}{!}{
          \begin{tabular}{|c|c|c|c|c|c|c|c|}
            \hline
            \rowcolor{teal!50}
            \textbf{MEMBRI} & \textbf{Responsabile} & \textbf{Amministratore} & \textbf{Analista} & \textbf{Progettista} & \textbf{Programmatore} & \textbf{Verificatore} & \textbf{Totale} \\ \hline
            Giulio Moretto & 1 & 1 & - & 6 & 13 & 3 & 24 \\ \hline
            Riccardo Simionato & 1 & 2 & - & 6 & 13 & 3 & 25 \\ \hline
            Marco Favaretto & - & 1 & - & 12 & 13 & 3 & 29 \\ \hline
            Marco Gobbo & - & - & - & 12 & 12 & 4 & 28 \\ \hline
            Matteo Camillo & 2 & 1 & - & 5 & 12 & 4 & 24 \\ \hline
            Edoardo Gallo & 2 & 2 & - & 3 & 12 & 4 & 23 \\ \hline
          \end{tabular}
        }
        \caption{Ore preventivate per la prima parte di progettazione e programmazione}
      \end{table}
    
    \subsubsection{Progettazione secondaria e codifica dei requisiti opzionali}
      \textbf{Periodo:} dal 15/04/2024 al 10/05/2024 
      \begin{table}[H]
        \centering
        \rowcolors{2}{black!15}{}
        \resizebox{\linewidth}{!}{
          \begin{tabular}{|c|c|c|c|c|c|c|c|}
            \hline
            \rowcolor{teal!50}
            \textbf{MEMBRI} & \textbf{Responsabile} & \textbf{Amministratore} & \textbf{Analista} & \textbf{Progettista} & \textbf{Programmatore} & \textbf{Verificatore} & \textbf{Totale} \\ \hline
            Giulio Moretto & 1 & 1 & - & 4 & 6 & 5 & 17 \\ \hline
            Riccardo Simionato & 1 & 2 & - & 3 & 6 & 3 & 15 \\ \hline
            Marco Favaretto & - & - & - & 6 & 6 & 4 & 16 \\ \hline
            Marco Gobbo & - & - & - & 6 & 6 & 4 & 16 \\ \hline
            Matteo Camillo & 2 & 1 & - & 3 & 6 & 5 & 17 \\ \hline
            Edoardo Gallo & 2 & 1 & - & 2 & 6 & 5 & 16 \\ \hline
          \end{tabular}
        }
        \caption{Ore preventivate per la realizzazione dei requisiti non obbligatori}
      \end{table}

    \subsubsection{Validazione finale e collaudo}
      \textbf{Periodo:} dal 10/05/2024 al 17/05/2024 
      \begin{table}[H]
        \centering
        \rowcolors{2}{black!15}{}
        \resizebox{\linewidth}{!}{
          \begin{tabular}{|c|c|c|c|c|c|c|c|}
            \hline
            \rowcolor{teal!50}
            \textbf{MEMBRI} & \textbf{Responsabile} & \textbf{Amministratore} & \textbf{Analista} & \textbf{Progettista} & \textbf{Programmatore} & \textbf{Verificatore} & \textbf{Totale} \\ \hline
            Giulio Moretto & 1 & 1 & - & - & - & 2 & 4 \\ \hline
            Riccardo Simionato & 1 & 2 & - & - & - & 2 & 5 \\ \hline
            Marco Favaretto & - & 1 & - & - & - & 2 & 3 \\ \hline
            Marco Gobbo & - & - & - & - & - & 3 & 3 \\ \hline
            Matteo Camillo & 2 & 1 & - & - & - & 2 & 5 \\ \hline
            Edoardo Gallo & 2 & 1 & - & - & - & 3 & 6 \\ \hline
          \end{tabular}
        }
        \caption{Ore preventivate per la fase finale della \textit{PB}}
      \end{table}
    
    \subsubsection{\textit{PB}}
      Qui sono presentate le ore e i costi preventivati per il periodo \textit{PB} complessivo.
      \begin{table}[H]
     	  \centering
     	  \rowcolors{2}{black!15}{}
     	  \resizebox{\linewidth}{!}{
     	    \begin{tabular}{|c|c|c|c|c|c|c|c|}
     			  \rowcolor{teal!50}
     			  \hline
   			    \textbf{MEMBRI} & \textbf{Responsabile} & \textbf{Amministratore} & \textbf{Analista} & \textbf{Progettista} &	\textbf{Programmatore} & \textbf{Verificatore} & \textbf{Totale} \\
   			    \hline
            Giulio Moretto & 4 & 5 & - & 10 & 19 & 13 & 51 \\ \hline
            Riccardo Simionato & 4 & 8 & - & 9 & 19 & 11 & 51 \\ \hline
            Marco Favaretto & - & 2 & - & 18 & 19 & 12 & 51 \\ \hline
            Marco Gobbo & - & - & - & 18 & 18 & 15 & 51 \\ \hline
            Matteo Camillo & 8 & 4 & - & 8 & 18 & 13 & 51 \\ \hline
            Edoardo	Gallo & 8 & 5 & - & 5 & 18 & 15 & 51 \\ \hline
          \end{tabular}
        }
        \caption{Piano ore preventivate per il secondo periodo}
      \end{table}

    \begin{table}[H]
    	\centering
    	\rowcolors{2}{black!15}{}
    	\resizebox{\linewidth}{!}{
    		\begin{tabular}{|c|c|c|c|c|c|c|c|}
    			\rowcolor{teal!50}
    			\hline
    			\textbf{MEMBRI} & \textbf{Responsabile} & \textbf{Amministratore} & \textbf{Analista} & \textbf{Progettista} &	\textbf{Programmatore} & \textbf{Verificatore} & \textbf{Totale} \\
    			\hline
          Totale per ruolo & 24 & 24 & 0 & 68 & 111 & 79 & 306 \\ \hline
          Costo per ora & \euro 30,00 & \euro 20,00 & \euro 25,00 & \euro 25,00 & \euro 15,00 & \euro 15,00 & - \\ \hline 
          Costo totale per ruolo & \euro 720,00 & \euro 480,00 & \euro - & \euro 1.700,00 & \euro 1.665,00 & \euro 1.185,00 & \euro 5.750,00 \\ \hline
        \end{tabular}
      }
      \caption{Tabella costi preventivati per il secondo periodo}
    \end{table}

\section{Consuntivo}
  Di seguito verranno mostrate le spese effettivamente avvenute nelle fasi citate
  \subsection{Periodo \textit{RTB}}
    La seguente tabella mostra le ore e i costi preventivati per questa fase di progetto per ogni ruolo, 
    affiancate alla differenza in termini di ore e i costi effettivi richiesti durante il primo periodo.
    \begin{table}[H]
      \centering
      \rowcolors{2}{black!15}{}
      \resizebox{\linewidth}{!}{
        \begin{tabular}{|c|c|c|c|c|c|c|c|}
          \hline
          \rowcolor{teal!50}
          \textbf{Ruolo} & \textbf{Ore previste} & \textbf{Differenza ore} & \textbf{Costo \euro/h} & \textbf{Costo preventivato} & \textbf{Costo effettivo} & \textbf{Differenza Costo} \\ \hline
          \textbf{Responsabile} & 24 & 0 & \euro 30,00 & \euro 720,00 & \euro 720,00 & \euro 0,00 \\ \hline
          \textbf{Amministratore} & 24 & -1,5 & \euro 20,00 & \euro 480,00 & \euro 510,00 & \euro 30,00 \\ \hline
          \textbf{Analista} & 72 & 0 & \euro 25,00 & \euro 1.800,00 & \euro 1.800,00 & \euro 0,00 \\ \hline
          \textbf{Progettista} & 40 & 14 & \euro 25,00 & \euro 1.000,00 & \euro 650,00 & - \euro 350,00 \\ \hline
          \textbf{Programmatore} & 21 & 1 & \euro 15,00 & \euro 315,00 & \euro 300,00 & - \euro 15,00 \\ \hline
          \textbf{Verificatore} & 65 & -1 & \euro 15,00 & \euro 975,00 & \euro 990,00 & \euro 15,00 \\ \hline \hline
          % doppia linea necessaria al fine di separare la prossima riga dal resto della tabella, senza necessariamente creare una nuova tabella
          \textbf{Totale} & 246 & 12,5	& - &	\euro 5.290,00 & \euro 4.970,00 & - \euro 320,00 \\ \hline
        \end{tabular}
      }
      \caption{Resoconto ore e costi \textit{RTB}}
    \end{table}
    Dal resoconto, si può notare come il gruppo abbia svolto questa prima fase entro i tempi previsti.
    Ciò è accaduto nonostante certi imprevisti e certe difficoltà, dovute principalmente ad impegni personali e all'inesperienza del \textit{team}. \\
    Ad esempio: nel mese di Marzo il \textit{team} ha cambiato piattaforma di \textit{ticketing} e gestione di \textit{spring}, passando da \textit{Trello} a \textit{Jira}.
    Ciò ha richiesto delle ore di lavoro per i ruoli di responsabile e amministratore aggiuntive a quelle preventivate.
    In seguito alla prima fase di revisione dell'\textit{RTB}, sono state necessarie ulteriori ore del ruolo di analista.
    Tuttavia, lo sviluppo del \textit{PoC} ha richeisto meno ore di progettazione, rispetto a quelle preventivate.


    Al netto di ciò, il consuntivo e il preventivo non presentano una notevole differenza. Abbiamo quindi ragione di credere che il gruppo continuerà 
    a lavorare in maniera simile anche nel prossimo periodo.

  \subsection{Periodo \textit{PB}}
    Le spese non sono presenti, in quanto il progetto non è ancora entrato in questa fase di sviluppo.


\end{document}
